% Zone „Das kennst du schon“ – aus bank/zinsrechnung/zone.jsonl (zusammenbau v0.1)
\einheitenkopf[zone][Kennst du schon]{Das kennst du schon}
\setcounter{aufgabe}{0}

%% TODO Titel: Ich-kann-Satz fehlt in der Bank (Platzhalter: Kettenname) – Nr. 1
%% TODO Anweisung: ein Satz über der ersten Teilaufgabe fehlt in der Bank – Nr. 1
\begin{aufgabe}{Prozentwert berechnen}
\begin{teile}
\teil 20\,\% von 40\,€ – wie viel Euro? \leerfeld[€]
\teil 3\,\% von 250\,€ – wie viel Euro? \leerfeld[€]
\end{teile}
\end{aufgabe}

%% TODO Titel: Ich-kann-Satz fehlt in der Bank (Platzhalter: Kettenname) – Nr. 2
%% TODO Anweisung: ein Satz über der ersten Teilaufgabe fehlt in der Bank – Nr. 2
\begin{aufgabe}{Prozentsatz berechnen}
\begin{teile}
\teil 5\,€ von 20\,€ – wie viel Prozent? \leerfeld[\%]
\teil 6,30\,€ Rabatt auf 90\,€ – wie viel Prozent? \leerfeld[\%]
\end{teile}
\end{aufgabe}

%% TODO Titel: Ich-kann-Satz fehlt in der Bank (Platzhalter: Kettenname) – Nr. 3
%% TODO Anweisung: ein Satz über der ersten Teilaufgabe fehlt in der Bank – Nr. 3
\begin{aufgabe}{Grundwert berechnen}
\begin{teile}
\teil 20\,\% sind 7\,€ – wie viel Euro ist das Ganze? \leerfeld[€]
\teil 3\,\% sind 4,50\,€ – wie viel Euro ist das Ganze? \leerfeld[€]
\end{teile}
\end{aufgabe}

%% TODO Titel: Ich-kann-Satz fehlt in der Bank (Platzhalter: Kettenname) – Nr. 4
%% TODO Anweisung: ein Satz über der ersten Teilaufgabe fehlt in der Bank – Nr. 4
\begin{aufgabe}{Dezimalzahlen addieren und multiplizieren, Euro-Beträge auf Cent runden (zwei Dezimalen), Ergebnis mit ≈}
\begin{teile}
\teil $2{,}50\,€ + 0{,}75\,€$ – wie viel Euro? \leerfeld[€]
\teil $345{,}60\,€ + 6{,}91\,€ + 50\,€$ – wie viel Euro? \leerfeld[€]
\end{teile}
\end{aufgabe}

%% TODO Titel: Ich-kann-Satz fehlt in der Bank (Platzhalter: Kettenname) – Nr. 5
%% TODO Anweisung: ein Satz über der ersten Teilaufgabe fehlt in der Bank – Nr. 5
\begin{aufgabe}{Dreisatz und fester Faktor („pro Monat“, „pro Tag“)}
\begin{teile}
\teil 3 Monate kosten 21\,€ – wie viel kosten 5 Monate? \leerfeld[€]
\teil Pro Tag 0,40\,€ – wie viel Euro für 25 Tage? \leerfeld[€]
\end{teile}
\end{aufgabe}

%% TODO Titel: Ich-kann-Satz fehlt in der Bank (Platzhalter: Kettenname) – Nr. 6
%% TODO Anweisung: ein Satz über der ersten Teilaufgabe fehlt in der Bank – Nr. 6
\begin{aufgabe}{Erhöhung um p \% als Wachstumsfaktor}
\begin{teile}
\teil „plus 50\,\%“ – mal welche Zahl? mal \leerfeld
\teil Ein Preis von 40\,€ steigt um 16\,\%. Rechne mit dem Faktor – neuer Preis? \leerfeld[€]
\end{teile}
\end{aufgabe}

%% TODO Titel: Ich-kann-Satz fehlt in der Bank (Platzhalter: Kettenname) – Nr. 7
%% TODO Anweisung: ein Satz über der ersten Teilaufgabe fehlt in der Bank – Nr. 7
\begin{aufgabe}{Potenz mit dem Taschenrechner}
\begin{teile}
\teil $2^{6}$ – welche Zahl? \leerfeld
\teil $1{,}04^{3}$ mit dem Taschenrechner, auf vier Stellen nach dem Komma gerundet? \leerfeld
\end{teile}
\end{aufgabe}

%% TODO Titel: Ich-kann-Satz fehlt in der Bank (Platzhalter: Kettenname) – Nr. 8
%% TODO Anweisung: ein Satz über der ersten Teilaufgabe fehlt in der Bank – Nr. 8
\begin{aufgabe}{Tabelle lesen und fortschreiben}
\begin{teile}
\teil Jede Woche kommen 4\,€ ins Sparschwein. Wie viel Euro sind es in Woche 4?

\sachtabelle{lr}{Woche & Betrag}{1 & 5\,€\\ 2 & 9\,€\\ 3 & 13\,€\\ 4 & \leerzelle}
\leerfeld[€]
\teil Neuer Stand = alter Stand − Ausgabe + Einnahme. Neuer Stand in Woche 2?

\sachtabelle{lrrrr}{Woche & alter Stand & Ausgabe & Einnahme & neuer Stand}{1 & 20,00\,€ & 3,40\,€ & 5,00\,€ & 21,60\,€\\ 2 & 21,60\,€ & 4,75\,€ & 5,00\,€ & \leerzelle}
\leerfeld[€]
\end{teile}
\end{aufgabe}

%% TODO Titel: Ich-kann-Satz fehlt in der Bank (Platzhalter: Kettenname) – Nr. 9
%% TODO Anweisung: ein Satz über der ersten Teilaufgabe fehlt in der Bank – Nr. 9
\begin{aufgabe}{Prozentwert berechnen – Fehler finden}
\begin{teile}
\teil 6\,\% von 650\,€. Ole rechnet: \rechnung{650 \cdot 0{,}6 &= 390} Antwort: 390\,€. Finde den Fehler und rechne richtig.
\end{teile}
\end{aufgabe}

%% TODO Titel: Ich-kann-Satz fehlt in der Bank (Platzhalter: Kettenname) – Nr. 10
%% TODO Anweisung: ein Satz über der ersten Teilaufgabe fehlt in der Bank – Nr. 10
\begin{aufgabe}{Prozentwert berechnen – selbst rechnen}
\begin{teile}
\teil 7\,\% von 300\,€ – wie viel Euro? \leerfeld[€]
\end{teile}
\end{aufgabe}
